\section{All finite level sets have dimension zero, and the complementary infinite level}\label{sec:finite-levels}

The sub-threshold result ($c<1/4$) uses the sharp classification. For
$c\ge1/4$ the level set $E_c$ is uncountable (Theorem~\ref{thm:gauge}), but we
show it still has Hausdorff dimension $0$. The entire covering argument ---
factorial growth of $Q_j$, uniform branching bound, surviving-center recursion,
periodicity reduction, Hausdorff-content estimate, and the countable-union
assembly --- is \emph{machine-checked at Layer~A}
(\texttt{EpiFormal/DimHCore.lean} + \texttt{EpiFormal/DimHFinite.lean}, base
axioms only, no \texttt{sorryAx}, and \emph{no} continued-fraction axioms).
The elementary proof below mirrors the formalization.

\begin{lemma}[Factorial growth, Layer~A in Lean]\label{lem:factorial}
For $j\ge1$, $Q_j\ge j!$. (Lean: \texttt{DimHCore.\allowbreak Qsel\_ge\_factorial}.)
\end{lemma}

\begin{proof}
We proceed by induction on $j\ge1$. For $j=1$, $Q_1=1=1!$. For $j=2$,
$Q_2=7\ge2=2!$. For $j\ge2$, every term $Q_i$ is positive, so the recurrence
\eqref{eq:rec} gives
\[
Q_{j+1}=(4j+2)Q_j+Q_{j-1}>(j+1)Q_j\ge(j+1)j!=(j+1)!,
\]
completing the induction.
\end{proof}

\begin{lemma}[Uniform branching bound, Layer~A in Lean]\label{lem:branch}
Fix $t\ge0$, $j\ge1$ and a parent center $c\in\Z$. The number of child centers
$k'\in\{0,1,\dots,Q_{j+1}\}$ with
\[
\left|\frac{c}{Q_j}-\frac{k'}{Q_{j+1}}\right|\le \frac{t}{jQ_j}+\frac{t}{(j+1)Q_{j+1}}
\]
is at most $B:=\lfloor16t\rfloor+1$, \emph{uniformly in $j$ and $c$}.
(Lean: \texttt{DimHCore.\allowbreak fiber\_card\_le}.)
\end{lemma}

\begin{proof}
Put $R_j:=t/(jQ_j)+t/((j+1)Q_{j+1})$. Multiplying the distance inequality
$|c/Q_j-k'/Q_{j+1}|\le R_j$ by $Q_{j+1}>0$ shows that every qualifying integer
$k'$ lies in the closed real interval $[A,A+W]$ with left endpoint
$A:=Q_{j+1}(c/Q_j-R_j)$ and width
\[
W:=2R_jQ_{j+1}=2t\left(\frac{Q_{j+1}}{jQ_j}+\frac{1}{j+1}\right).
\]
Since $Q_0=Q_1=1$ and $Q_{i-1}\le Q_i$ for all $i\ge1$, the recurrence
\eqref{eq:rec} implies $Q_{j+1}=(4j+2)Q_j+Q_{j-1}\le(4j+3)Q_j\le7jQ_j$ for all
$j\ge1$, while $1/(j+1)\le1$. Hence
\[
W\le 2t(7+1)=16t.
\]
For any qualifying $k'$, we have $0\le k'-A\le W\le16t$, so the map
$k'\mapsto\lfloor k'-A\rfloor$ takes values in $\{0,1,\dots,\lfloor16t\rfloor\}$
and is injective on integers (if $k'_1<k'_2$, then $k'_2-A\ge(k'_1-A)+1$, so
$\lfloor k'_2-A\rfloor\ge\lfloor k'_1-A\rfloor+1$). Therefore the number of
qualifying integers $k'$ is at most $\lfloor16t\rfloor+1=B$.
\end{proof}

\begin{theorem}[Finite level sets, Layer~A in Lean]\label{thm:finite}
For every $c\in\R$, $\dimH E_c=0$. Consequently
$F:=\{x\in\R:\Lam(x)<\infty\}=\bigcup_{n\ge1}E_n$ has $\dimH F=0$.
(Lean: \texttt{DimHFinite.\allowbreak dimH\_levelSet\_finite}; \texttt{\#print axioms}
returns \texttt{[propext, Classical.choice, Quot.sound]} only.)
\end{theorem}

\begin{proof}
Fix $c\in\R$ and choose a rational or real threshold $t>\max(c,0)$ (in Lean,
$t=\max(c,0)+1>0$). First note that for every $m\in\Z$ and $j\ge1$, since
$Q_j\in\Z$, we have $\|Q_j(x+m)\|=\|Q_jx+Q_jm\|=\|Q_jx\|$, and hence
$\Lam(x+m)=\Lam(x)$. Thus $E_c$ is invariant under integer translations and
decomposes as the countable union
\[
E_c=\bigcup_{m\in\Z}\bigl(m+(E_c\cap[0,1])\bigr).
\]
Because integer translations are isometries and Hausdorff dimension is stable
under countable unions ($\dimH(\bigcup_m X_m)=\sup_m\dimH X_m$), it suffices to
prove $\dimH(E_c\cap[0,1])=0$.

For each $j\ge1$, define $r_j:=t/(jQ_j)$ and
$A_j(t):=\{x\in\R:\|Q_jx\|\le t/j\}$. If $x\in E_c\cap[0,1]$, then
$\limsup_{j\to\infty}j\|Q_jx\|=\Lam(x)\le c<t$, which by definition of the limit
superior means that $j\|Q_jx\|\le t$ \emph{eventually} (a liminf-type tail
condition, \emph{not} a limsup set). Therefore there exists an integer $J\ge1$
such that $x\in A_j(t)$ for all $j\ge J$. Consequently,
\[
E_c\cap[0,1]\subseteq\bigcup_{J=1}^\infty C_J,\qquad
C_J:=[0,1]\cap\bigcap_{j=J}^\infty A_j(t).
\]
By countable stability again, it suffices to fix $J\ge1$ and prove that
$\dimH C_J=0$, i.e.\ $\mathcal H^s(C_J)=0$ for every $s>0$.

For $j\ge J$ and $k\in\Z$, let $I_j(k):=[k/Q_j-r_j,k/Q_j+r_j]$ be the closed
interval of center $k/Q_j$ and radius $r_j$ (so $\operatorname{diam}I_j(k)=2r_j$).
For any $x\in C_J$ and $j\ge J$, let $k_j(x):=\mathrm{round}(Q_jx)\in\Z$ be a
nearest integer to $Q_jx$. Since $x\in[0,1]$ and $Q_j\ge1$, we have
$0\le Q_jx\le Q_j$, so $k_j(x)\in\{0,1,\dots,Q_j\}$; moreover, $x\in A_j(t)$
gives $|Q_jx-k_j(x)|=\|Q_jx\|\le t/j$, and dividing by $Q_j$ yields
$|x-k_j(x)/Q_j|\le r_j$, that is, $x\in I_j(k_j(x))$. Furthermore, at two
consecutive levels $j$ and $j+1$, the same point $x\in C_J$ belongs to both
$I_j(k_j(x))$ and $I_{j+1}(k_{j+1}(x))$, so the triangle inequality gives the
parent--child center compatibility condition
\begin{equation}\label{eq:parent-child-tri}
\left|\frac{k_j(x)}{Q_j}-\frac{k_{j+1}(x)}{Q_{j+1}}\right|
\le\left|\frac{k_j(x)}{Q_j}-x\right|+\left|x-\frac{k_{j+1}(x)}{Q_{j+1}}\right|
\le r_j+r_{j+1}=\frac{t}{jQ_j}+\frac{t}{(j+1)Q_{j+1}}.
\end{equation}

Define the surviving center sets $S_n\subseteq\{0,1,\dots,Q_{J+n}\}$ for
$n\ge0$ inductively by setting $S_0:=\{0,1,\dots,Q_J\}$ and, for $n\ge0$,
letting $S_{n+1}$ be the set of all $k'\in\{0,1,\dots,Q_{J+n+1}\}$ for which
there exists $k\in S_n$ satisfying
$|k/Q_{J+n}-k'/Q_{J+n+1}|\le r_{J+n}+r_{J+n+1}$. By induction on $n\ge0$ using
\eqref{eq:parent-child-tri}, every $x\in C_J$ satisfies $k_{J+n}(x)\in S_n$ for
all $n\ge0$, and hence
\[
C_J\subseteq\bigcup_{k\in S_n} I_{J+n}(k)\qquad\text{for all }n\ge0.
\]
By Lemma~\ref{lem:branch}, for each parent $k\in S_n$ the number of children
$k'\in S_{n+1}$ is at most $B=\lfloor16t\rfloor+1$. Since $|S_0|=Q_J+1$,
induction gives $|S_n|\le(Q_J+1)B^n$ for all $n\ge0$. Each interval
$I_{J+n}(k)$ has diameter
\[
\operatorname{diam}I_{J+n}(k)=2r_{J+n}=\frac{2t}{(J+n)Q_{J+n}}
\le\frac{2t}{(J+n)!}\longrightarrow0\qquad(n\to\infty),
\]
where we used $Q_{J+n}\ge(J+n)!$ from Lemma~\ref{lem:factorial} and $J+n\ge1$.

Now fix any $s>0$ and any gauge scale $\delta>0$. For all sufficiently large
$n$, the intervals $\{I_{J+n}(k):k\in S_n\}$ form a finite $\delta$-cover of
$C_J$, and by Lemma~\ref{lem:factorial} ($Q_{J+n}\ge(J+n)!$) their $s$-dimensional
Hausdorff covering cost satisfies
\[
\sum_{k\in S_n}\bigl(\operatorname{diam}I_{J+n}(k)\bigr)^s
\le |S_n|\left(\frac{2t}{(J+n)Q_{J+n}}\right)^s
\le(Q_J+1)(2t)^s B^{-J}\left(\frac{(B^{1/s})^{J+n}}{(J+n)!}\right)^s.
\]
Using the elementary factorial lower bound $m!\ge(m/\ee)^m$ for $m=J+n\ge1$
(or the convergent ratio $(B^{1/s})^{m+1}/(m+1)!\div((B^{1/s})^m/m!)=B^{1/s}/(m+1)<1/2$
for $m+1>2B^{1/s}$, formalized in mathlib as
\texttt{FloorSemiring.tendsto\_pow\_div\_factorial\_atTop}), we have
\[
0\le\frac{(B^{1/s})^m}{m!}\le\left(\frac{\ee B^{1/s}}{m}\right)^m\longrightarrow0
\qquad(m=J+n\to\infty),
\]
and raising to the positive power $s>0$ shows that the covering cost tends to
$0$ as $n\to\infty$. Thus $\mathcal H_\delta^s(C_J)=0$ for every $\delta>0$,
so $\mathcal H^s(C_J)=\lim_{\delta\downarrow0}\mathcal H_\delta^s(C_J)=0$ for
every $s>0$, which gives $\dimH C_J=0$. Taking the countable union over
$J\ge1$ and $m\in\Z$ yields $\dimH E_c=0$. Finally,
$F=\{x\in\R:\Lam(x)<\infty\}=\bigcup_{n=1}^\infty E_n$ is a countable union of
sets of Hausdorff dimension $0$, so $\dimH F=0$.
\end{proof}

\begin{corollary}[The complementary infinite level, Layer~A in Lean]\label{cor:infinite-level}
The set $F=\{x\in\R:\Lam(x)<\infty\}$ has Lebesgue measure zero, and its
complement $F_\infty:=\{x\in\R:\Lam(x)=\infty\}=\R\setminus F$ has full
Lebesgue measure on $\R$. Moreover, for every nondegenerate interval
$[a,b]\subset\R$ ($a<b$), one has
$\operatorname{Leb}(F_\infty\cap[a,b])=b-a>0$ and
\[
\dimH(F_\infty\cap[a,b])=\dimH F_\infty=1.
\]
(Lean, module \texttt{EpiFormal/\allowbreak GaugeExact.lean}: Layer~A, base axioms only.)
\end{corollary}

\begin{proof}
That $\operatorname{Leb}(F)=0$ follows in two independent ways. First, by
Theorem~\ref{thm:finite}, $\dimH F=0<1$, which implies $\mathcal H^1(F)=0$;
since one-dimensional Hausdorff measure $\mathcal H^1$ on $\R$ coincides with
outer Lebesgue measure $\operatorname{Leb}$, we obtain $\operatorname{Leb}(F)=0$.

Second, one can see $\operatorname{Leb}(F)=0$ directly by an elementary
measure-preserving calculation without invoking Hausdorff dimension. For an
integer $m\ge1$ and $j>2m$, the circle endomorphism $T_{Q_j}:\R/\Z\to\R/\Z$,
$x\mapsto Q_jx\bmod1$, preserves normalized Lebesgue measure on $\R/\Z$ (each
arc of length $\ell$ has $Q_j$ disjoint preimage arcs of length $\ell/Q_j$).
The set $\bar A_j(m):=\{x\in\R/\Z:\|Q_jx\|\le m/j\}=T_{Q_j}^{-1}([-m/j,m/j]\bmod1)$
therefore has circle Lebesgue measure $2m/j$. For any fixed starting index
$J\ge1$ and every $j\ge\max(J,2m+1)$, monotonic inclusion gives
\[
\operatorname{Leb}_{\R/\Z}\left(\bigcap_{i=J}^\infty \bar A_i(m)\right)
\le\operatorname{Leb}_{\R/\Z}\bigl(\bar A_j(m)\bigr)=\frac{2m}{j}\longrightarrow0
\qquad(j\to\infty),
\]
so $\operatorname{Leb}_{\R/\Z}(\bigcap_{i\ge J}\bar A_i(m))=0$. Every point
$x\in\R/\Z$ with $\Lam(x)<\infty$ belongs to $\bigcap_{i\ge J}\bar A_i(m)$ for
some integers $m>\Lam(x)$ and $J\ge1$. Countable subadditivity over $m\ge1$ and
$J\ge1$ yields $\operatorname{Leb}_{\R/\Z}(\{\Lam<\infty\})=0$ on the circle,
and countable union over integer translates $[k,k+1]$ ($k\in\Z$) yields
$\operatorname{Leb}(F)=0$ on $\R$.

Now let $[a,b]\subset\R$ be any nondegenerate closed interval ($a<b$). Since
$[a,b]=(F_\infty\cap[a,b])\cup(F\cap[a,b])$ and $\operatorname{Leb}(F\cap[a,b])\le\operatorname{Leb}(F)=0$,
additivity of Lebesgue measure gives
\[
\mathcal H^1(F_\infty\cap[a,b])=\operatorname{Leb}(F_\infty\cap[a,b])
=\operatorname{Leb}([a,b])-\operatorname{Leb}(F\cap[a,b])=b-a>0.
\]
By the definition of Hausdorff dimension, $\mathcal H^1(F_\infty\cap[a,b])>0$
implies $\dimH(F_\infty\cap[a,b])\ge1$. On the other hand, any subset of $\R$
has Hausdorff dimension at most $\dimH\R=1$, so
\[
1\le\dimH(F_\infty\cap[a,b])\le\dimH F_\infty\le\dimH\R=1.
\]
Hence $\dimH(F_\infty\cap[a,b])=\dimH F_\infty=1$. Together with
Theorem~\ref{thm:finite}, this establishes the sharp metric dichotomy: $\Lam$
is infinite Lebesgue-almost everywhere (with local Hausdorff dimension~$1$ on
every interval), while the set where $\Lam$ is finite has Hausdorff
dimension~$0$.
\end{proof}

\paragraph{Lean status of the assembly.} All steps of Theorem~\ref{thm:finite}
are kernel-checked at Layer~A. The surviving centers are encoded as a
\texttt{Finset} recursion (\texttt{survSet}), not as paths:
$S_0=\{0,\dots,Q_J\}$ and $S_{n+1}=\bigcup_{k\in S_n}\mathrm{branchSet}(k)$,
with $|S_n|\le(Q_J+1)B^n$ by \texttt{Finset.card\_biUnion\_le\_card\_mul} and
Lemma~\ref{lem:branch} (\texttt{card\_survSet\_le}). Membership of the center
sequence $\mathrm{round}(Q_{J+n}x)$ in $S_n$ is proved by induction via the
triangle inequality (\texttt{centerOf\_mem\_survSet}), giving the cover
(\texttt{tailSet\_cover}). The content bound
$|S_n|(2t/((J+n)Q_{J+n}))^s\to0$ uses $Q_j\ge j!$ and
\texttt{FloorSemiring.\allowbreak tendsto\_pow\_div\_factorial\_atTop} with
$c=B^{1/s}$ (\texttt{tendsto\_contentBound}); no asymptotic notation is used.
Feeding this into \texttt{hausdorffMeasure\_le\_liminf\_sum} yields
$\mu H^s=0$ for all $s>0$ (\texttt{hausdorffMeasure\_tailSet\_zero}), hence
$\dimH=0$ (\texttt{dimH\_tailSet\_zero}). Periodicity
$\|Q_j(x+m)\|=\|Q_jx\|$ (\texttt{apprSeq\_add\_intCast}) reduces general $x$ to
$[0,1]$; integer shifts are $1$-Lipschitz and countable unions preserve
dimension $0$ (\texttt{dimH\_iUnion}). The sub-threshold case $c<1/4$ is
\emph{not} used: a monotonicity reduction to $c'=\max(c,0)$ lets the covering
argument handle all $c$ uniformly, so the theorem needs no continued-fraction
axioms.

